\Needspace{18\baselineskip}
\subsection{The least eigenvalue of Lee curvature}
\label{subsec:lck-minimum-lee-curvature}

Fix a strictly LCK metric $\omega$ with Lee form $\lambda$.
The volume along $\omega+t\kappa_\lambda$ increases as long
as this form is positive. If the least relative eigenvalue of
$\kappa_\lambda$ is a negative constant, the first boundary
form has a kernel everywhere and its total volume is zero.
This contradiction gives a pointwise positivity criterion.

Set $Q=\lambda\wedge\lambda^{\mathrm c}\geq0$.
The form $Q$ is nonzero on an open set, since $\lambda$ is
not exact and hence is not identically zero.

\Needspace{14\baselineskip}
\begin{proposition}\label{lck:lee-volume-ray}\index{Lee form!volume polynomial}
On the interval where $\eta_t=\omega+t\kappa_\lambda>0$, the
polynomial $V(t)=\int_X\eta_t^n$ has degree at most $n-1$ and
\[
 V'(t)=n(n-1)\int_XQ\wedge\omega\wedge\eta_t^{n-2}>0.
\]
\end{proposition}
\begin{proof}
Since $\mathrm d\eta_t=\lambda\wedge\omega$ and
$\mathrm d\lambda^{\mathrm c}=-\kappa_\lambda$, Stokes' theorem gives
\[
 0=\int_X\mathrm d(\lambda^{\mathrm c}\wedge\eta_t^{n-1})
  =-\int_X\kappa_\lambda\wedge\eta_t^{n-1}
    +(n-1)\int_XQ\wedge\omega\wedge\eta_t^{n-2}.
\]
Differentiating the wedge polynomial proves the formula.
The integrand on the right is positive wherever $\lambda\ne0$,
including when $n=2$. Its integral is therefore strictly positive.
Finally, $\kappa_\lambda$ is closed and de Rham exact, so
\[
 \int_X\kappa_\lambda^n
 =-\int_X\mathrm d(\lambda^{\mathrm c}\wedge
                         \kappa_\lambda^{n-1})=0.
\]
This removes the coefficient of degree $n$.
\end{proof}

\Needspace{16\baselineskip}
\begin{theorem}\label{lck:constant-minimum-lee}\index{Lee form!constant least curvature eigenvalue}
If the least eigenvalue of $\kappa_\lambda$ relative to
$\omega$ is constant on $X$, then it is zero and
$\kappa_\lambda$ is semipositive. Its ordinary Bott--Chern
class is nonzero and pseudoeffective.

If all relative eigenvalues are constant, the rank is constant
and lies between $1$ and $n-1$. In particular, the conclusion
holds on a compact quotient $\Gamma\backslash G$ with
left-invariant complex structure and strictly LCK tensors.
\end{theorem}
\begin{proof}
Let $r$ be the constant least relative eigenvalue. If $r<0$,
put $t_0=-1/r$. Then $\eta_t>0$ for $0\leq t<t_0$,
whereas $\eta_{t_0}$ is semipositive with a kernel at every
point. Thus $V(t_0)=0$. The preceding proposition and
polynomial continuity instead give $V(t_0)\geq V(0)>0$.
This proves $r\geq0$ and semipositivity.

The nonnegative smooth form $\kappa_\lambda^n$ has integral
zero by the same proposition, hence vanishes everywhere.
Thus the least eigenvalue is zero at every point. Applying
the Lee degree \eqref{eq:lck-lee-degree} to a conformal
Gauduchon representative, as above, proves that the class is
nonzero. The form $\kappa_\lambda$ itself represents it positively.

Constant spectrum gives constant rank. The positive pairing
\eqref{eq:lck-lee-degree} excludes rank zero, and the vanishing
top power excludes rank $n$. On the stated compact quotient,
the relative matrix is constant in a left-invariant frame,
which proves the last assertion.
\end{proof}

The class is also nef: the same representative is bounded below by
$-\epsilon h$ for every Hermitian form $h$ and every $\epsilon>0$.
Thus \cref{prop:lck-balanced-obstruction} excludes balanced metrics
under the least-eigenvalue hypothesis.

The hypothesis concerns the chosen LCK metric. The other
eigenvalues may vary, as the following examples show. They also
explain why first degeneration of a volume ray need not force
its total volume to vanish without the constant-minimum condition.

\Needspace{16\baselineskip}
\begin{example}\label{lck:indefinite-hopf-lee}\index{Hopf manifold!indefinite Lee curvature}
On the scalar Hopf manifold
$H_n=(\mathbb C^n\setminus\{0\})/\langle z\mapsto z/2\rangle$,
for every $n\geq2$, there is a strictly LCK metric with
indefinite Lee curvature and a smooth semipositive representative
of its ordinary Bott--Chern class. The first boundary reached
for increasing $t>0$ in $\omega+t\kappa_\lambda$ has positive
total volume.
\end{example}
\begin{proof}
Fix $0<\epsilon<1$ and put
\[
 \rho=\sum_j|z_j|^2,\qquad
 \Phi=\rho+\epsilon\operatorname{Re}(z_1^2),\qquad
 \omega=\frac{\ddc\rho}{\Phi},\qquad
 \lambda=-\mathrm d\log\Phi.
\]
The bounds $(1-\epsilon)\rho\leq\Phi\leq(1+\epsilon)\rho$
give positivity. Since $\ddc\Phi=\ddc\rho$ and both $\Phi$
and $\ddc\rho$ scale by $1/4$ under $z\mapsto z/2$,
the metric and Lee form descend and satisfy the LCK equation.
Polar coordinates identify the quotient with
$S^{2n-1}\times(\mathbb R/(\log2)\mathbb Z)$, so it is compact
and connected. The Lee period from $z$ to $z/2$ is $2\log2$,
which proves strictness.

Here $\kappa_\lambda=\ddc\log\Phi$. With
$F_i=\bar z_i+\epsilon z_1\delta_{i1}$, its complex Hessian is
\[
 (\log\Phi)_{i\bar j}
 =\frac{\delta_{ij}}{\Phi}
       -\frac{F_i\overline{F_j}}{\Phi^2}.
\]
Relative to $\omega$, the eigenvalues are $1$ with multiplicity
$n-1$ and $h=1-|F|^2/\Phi$. Writing $z_1=x+\mathrm i y$ and
$B=\sum_{j>1}|z_j|^2$, one finds
\[
 h=\frac{\epsilon((1-\epsilon)y^2-(1+\epsilon)x^2)}
           {(1+\epsilon)x^2+(1-\epsilon)y^2+B},
 \qquad -\epsilon\leq h\leq\epsilon.
\]
At $(1,0,\ldots,0)$ the spectrum is
$-\epsilon,1,\ldots,1$, hence is indefinite. The bounded smooth
function $u=\log\rho-\log\Phi$ is deck invariant, and
\[
 \kappa_\lambda+\ddc u=\ddc\log\rho\geq0,
\]
by the Cauchy--Schwarz inequality for the Hessian of $\log\rho$.
The correction also gives the standard Hopf conformal representative:
$\mathrm e^{-u}\omega=\ddc\rho/\rho$.

The first boundary occurs at $t_0=1/\epsilon$. At that value,
\[
 1+h/\epsilon=\frac{2(1-\epsilon)y^2+B}{\Phi}.
\]
It vanishes only on the real $z_1$ rays with the remaining
coordinates zero. The other relative eigenvalues of
$\omega+t_0\kappa_\lambda$ equal $1+1/\epsilon>0$.
The boundary form is therefore semipositive and its top power
is positive on a nonempty open set, proving positive total volume.
\end{proof}

\Needspace{16\baselineskip}
\begin{example}\label{lck:variable-positive-lee-spectrum}\index{Hopf manifold!variable Lee spectrum}
For every $n\geq2$, the same $H_n$ carries a family of strictly
LCK metrics whose least Lee-curvature eigenvalue is identically
zero and whose remaining relative spectrum is nonconstant.
Thus constancy of the least eigenvalue is strictly weaker than
constancy of the whole spectrum.
\end{example}
\begin{proof}
Put $q=|z_1|^2$ and, for $0<\epsilon<1$, define
\[
 \Psi=\rho+\epsilon\frac{q^2}{\rho},\qquad
 \omega_\epsilon=\frac{\ddc\Psi}{\rho},\qquad
 \lambda=-\mathrm d\log\rho.
\]
For $a=q/\rho\in[0,1]$, direct differentiation gives
\[
 \left(\frac{q^2}{\rho}\right)_{i\bar j}
 =\frac{2}{\rho}(q_i-a\rho_i)(q_{\bar j}-a\rho_{\bar j})
       +2a\delta_{i1}\delta_{j1}-a^2\delta_{ij}.
\]
The first two terms are positive semidefinite and $a^2\leq1$,
so $\ddc\Psi\geq(1-\epsilon)\ddc\rho>0$.
The potential is smooth on the punctured space and scales by
$1/4$ under the deck map. Hence $\omega_\epsilon$ descends,
satisfies the LCK equation and has Lee period $2\log2$.

Its Lee curvature is $\ddc\log\rho$, which is semipositive
of rank $n-1$. The least relative eigenvalue is consequently
zero everywhere. At $(1,0,\ldots,0)$, the Hessian of
$q^2/\rho$ is $\operatorname{diag}(1,-1,\ldots,-1)$, giving
the spectrum
\[
 0,(1-\epsilon)^{-1},\ldots,(1-\epsilon)^{-1}.
\]
At $(0,1,0,\ldots,0)$ the term $q^2/\rho$ vanishes to fourth
order, so its Hessian is zero and the spectrum is
$0,1,\ldots,1$. These spectra differ for every $\epsilon>0$.
\end{proof}
